% =====================================================================
\chapter{Inducción simple y fuerte; buen orden}
\label{tema:5}
% =====================================================================

Muchos enunciados del curso dependen de un número natural: la longitud de una secuencia, el número de capas de una red, el paso de un algoritmo iterativo. La inducción matemática permite demostrarlos para todos los valores a la vez, y es la técnica de demostración más usada en informática. Este tema presenta la inducción simple, la inducción fuerte y el principio del buen orden, y muestra que son tres formas de la misma idea. Se siguen el capítulo~10 de~\cite{hammack}, los capítulos~2 y~5 de~\cite{mcs} y las §6.1, 6.2 y 6.4 de~\cite{velleman}.

\section{El principio de inducción}

\begin{teorema}[Principio de inducción]\label{teo:induccion}
Sea $P(n)$ un predicado sobre $\N_0$. Si
\begin{enumerate}
  \item (\emph{caso base}) $P(0)$ es verdadera, y
  \item (\emph{paso de inducción}) para todo $n\in\N_0$, $P(n)\Rightarrow P(n+1)$,
\end{enumerate}
entonces $P(n)$ es verdadera para todo $n\in\N_0$.
\end{teorema}

La imagen habitual es la de una fila infinita de fichas de dominó: la primera cae, y cada una que cae derriba a la siguiente. En el paso de inducción, la suposición $P(n)$ se llama \emph{hipótesis de inducción}; no se está suponiendo lo que se quiere demostrar, sino probando un condicional para un $n$ arbitrario. Si el caso base es $P(n_0)$ con $n_0\in\Z$, la conclusión vale para todo $n\ge n_0$. El \cref{teo:induccion} se toma como axioma de los números naturales o se deduce del principio del buen orden (\cref{prop:bo-induccion}).

\begin{ejemplo}[Suma de los primeros naturales]\label{ej:gauss}
Se demuestra que, para todo $n\in\N$, $1+2+\dots+n=\dfrac{n(n+1)}{2}$.

\emph{Demostración.} Sea $P(n)$ la igualdad del enunciado. \emph{Caso base:} para $n=1$, el miembro izquierdo vale~1 y el derecho $1\cdot 2/2=1$. \emph{Paso de inducción:} sea $n\ge 1$ y supóngase $P(n)$. Entonces
\[
1+\dots+n+(n+1)\overset{P(n)}{=}\frac{n(n+1)}{2}+(n+1)=\frac{(n+1)(n+2)}{2},
\]
que es $P(n+1)$. Por el principio de inducción, $P(n)$ vale para todo $n\in\N$. \hfill$\blacksquare$
\end{ejemplo}

\begin{ejemplo}[Serie geométrica finita]\label{ej:geometrica}
Para $r\in\R$, $r\neq 1$, y $n\in\N_0$:
\[
\sum_{k=0}^{n}r^{k}=\frac{r^{n+1}-1}{r-1}.
\]
\emph{Caso base:} $n=0$ da $1=(r-1)/(r-1)$. \emph{Paso:} si vale para $n$, entonces
$\sum_{k=0}^{n+1}r^{k}=\frac{r^{n+1}-1}{r-1}+r^{n+1}=\frac{r^{n+1}-1+r^{n+2}-r^{n+1}}{r-1}=\frac{r^{n+2}-1}{r-1}$. \hfill$\blacksquare$
\end{ejemplo}

\begin{ejemplo}[Una desigualdad]\label{ej:desig-induccion}
Para todo $n\in\N_0$, $2^{n}\ge n+1$. \emph{Caso base:} $2^{0}=1\ge 1$. \emph{Paso:} si $2^{n}\ge n+1$, entonces $2^{n+1}=2\cdot 2^{n}\ge 2(n+1)=n+2+n\ge n+2$. \hfill$\blacksquare$
\end{ejemplo}

En las desigualdades, el paso de inducción suele necesitar una estimación intermedia, como aquí $2n+2\ge n+2$. Conviene escribir primero qué hay que obtener, $P(n+1)$, y buscar cómo llegar desde $P(n)$.

\begin{ejemplo}[Divisibilidad]\label{ej:div-induccion}
Para todo $n\in\N_0$, $3\divides 4^{n}-1$. \emph{Caso base:} $4^{0}-1=0=3\cdot 0$. \emph{Paso:} si $4^{n}-1=3m$, entonces $4^{n+1}-1=4\cdot 4^{n}-1=4(3m+1)-1=3(4m+1)$. \hfill$\blacksquare$
\end{ejemplo}

\section{Inducción fuerte}

\begin{teorema}[Inducción fuerte]\label{teo:induccion-fuerte}
Sea $P(n)$ un predicado sobre los enteros $n\ge n_0$. Si $P(n_0)$ es verdadera y, para todo $n\ge n_0$,
\[
\bigl(P(n_0)\wedge P(n_0+1)\wedge\dots\wedge P(n)\bigr)\Rightarrow P(n+1),
\]
entonces $P(n)$ es verdadera para todo $n\ge n_0$.
\end{teorema}

\begin{proof}
Sea $Q(n)$ el predicado «$P(n_0),\dots,P(n)$ son verdaderas». $Q(n_0)$ es $P(n_0)$, verdadera. Si $Q(n)$ es verdadera, la hipótesis da $P(n+1)$, y por tanto $Q(n+1)$. Por inducción simple, $Q(n)$ vale para todo $n\ge n_0$, y en particular $P(n)$.
\end{proof}

La inducción fuerte es útil cuando $P(n+1)$ depende no de $P(n)$, sino de casos anteriores más lejanos, como al descomponer un número en factores o en las recurrencias de varios términos.

\begin{teorema}[Existencia de la factorización en primos]\label{teo:factorizacion}
Todo entero $n\ge 2$ es primo o producto de primos.
\end{teorema}

\begin{proof}
Por inducción fuerte sobre $n$. \emph{Caso base:} $2$ es primo. \emph{Paso:} sea $n\ge 2$ y supóngase el enunciado cierto para $2,3,\dots,n$. Si $n+1$ es primo, no hay nada que probar. Si no lo es, $n+1=ab$ con $2\le a,b\le n$. Por la hipótesis de inducción, $a$ y $b$ son primos o productos de primos, y entonces $n+1=ab$ es producto de primos.
\end{proof}

\begin{ejemplo}[Varios casos base]\label{ej:sellos}
Todo entero $n\ge 8$ se escribe como $n=3a+5b$ con $a,b\in\N_0$. \emph{Casos base:} $8=3+5$, $9=3\cdot 3$, $10=5\cdot 2$. \emph{Paso:} sea $n\ge 10$ y supóngase el enunciado para $8,\dots,n$. Entonces $n-2\ge 8$, luego $n-2=3a+5b$ y $n+1=3(a+1)+5b$. \hfill$\blacksquare$

Se necesitan tres casos base porque el paso usa el caso $n-2$, que debe ser al menos~8; con un solo caso base, el argumento no cubriría $9$ ni $10$.
\end{ejemplo}

\section{El principio del buen orden}

\begin{teorema}[Principio del buen orden]\label{teo:buen-orden}
Todo subconjunto no vacío de $\N_0$ tiene un elemento mínimo.
\end{teorema}

\begin{proof}
Sea $S\subseteq\N_0$ sin elemento mínimo; se demuestra que $S=\varnothing$. Sea $P(n)$: «ninguno de $0,1,\dots,n$ pertenece a $S$». $P(0)$ es cierta: si $0\in S$, sería su mínimo. Si $P(n)$ es cierta y $n+1\in S$, entonces $n+1$ sería el mínimo de $S$, porque ningún número menor está en $S$; luego $n+1\notin S$ y vale $P(n+1)$. Por inducción, $P(n)$ vale para todo $n$, y ningún natural está en $S$.
\end{proof}

\begin{proposicion}\label{prop:bo-induccion}
El principio del buen orden implica el principio de inducción.
\end{proposicion}

\begin{proof}
Sea $P$ con $P(0)$ y $P(n)\Rightarrow P(n+1)$ para todo $n$. Supóngase, por reducción al absurdo, que el conjunto $C=\{n\in\N_0:\neg P(n)\}$ de contraejemplos es no vacío, y sea $m$ su mínimo. Como $P(0)$ es cierta, $m\ge 1$, y $m-1\in\N_0$ no está en $C$ por ser menor que $m$; luego $P(m-1)$ es cierta. El paso de inducción da $P(m)$, en contra de $m\in C$.
\end{proof}

La demostración anterior es el modelo de las \emph{pruebas por el menor contraejemplo}~\cite[§2.2]{mcs}, \cite[§10.3]{hammack}: para demostrar $\forall n\,P(n)$, se supone que el conjunto de contraejemplos no es vacío, se toma el menor y se construye un contraejemplo aún menor, o se llega a otra contradicción.

\begin{proposicion}[Existencia en la división euclídea]\label{prop:division-existencia}
Sean $a\in\Z$ y $b\in\N$. Existen $q,r\in\Z$ con $a=qb+r$ y $0\le r<b$.
\end{proposicion}

\begin{proof}
Sea $S=\{a-qb: q\in\Z,\ a-qb\ge 0\}\subseteq\N_0$. Es no vacío: con $q=-\abs{a}$ se tiene $a+\abs{a}b\ge a+\abs{a}\ge 0$, porque $b\ge 1$. Por el buen orden tiene un mínimo $r=a-qb\ge 0$. Si fuera $r\ge b$, entonces $r-b=a-(q+1)b\ge 0$ estaría en $S$ y sería menor que $r$, en contra de la minimalidad. Luego $0\le r<b$.
\end{proof}

\section{Definiciones recursivas}

Los sumatorios, los productorios, las potencias y muchas sucesiones se definen \emph{recursivamente}: se da el primer valor y una regla para obtener cada valor a partir de los anteriores. Por ejemplo, $\sum_{k=1}^{0}a_k=0$ y $\sum_{k=1}^{n+1}a_k=\sum_{k=1}^{n}a_k+a_{n+1}$; la sucesión de Fibonacci se define por $F_1=F_2=1$ y $F_{n+2}=F_{n+1}+F_n$. La inducción es la forma natural de demostrar propiedades de objetos definidos así: el caso base comprueba la definición inicial y el paso usa la regla recursiva~\cite[§10.5]{hammack}.

\begin{observacion}[Errores frecuentes]\label{obs:5-errores}
(a) Omitir el caso base. Con $P(n)$: «$n=n+1$», el paso $P(n)\Rightarrow P(n+1)$ es correcto (sumando~1), pero $P(0)$ es falsa y el enunciado también.
(b) Un paso que no vale para todo $n$. La célebre «demostración» de que todos los caballos tienen el mismo color usa que dos grupos de $n$ caballos se solapan, lo que falla al pasar de $n=1$ a $n=2$ (\cref{eje:5-caballos}).
(c) Usar $P(n+1)$ dentro del paso: hay que llegar a $P(n+1)$, no partir de ella.
(d) En inducción fuerte, no incluir todos los casos base que el paso necesita (\cref{ej:sellos}).
\end{observacion}

\begin{observacion}[Conexión con el curso]
La inducción aparece en todo el curso. En el módulo~1, la regla de la cadena de la probabilidad se demuestra por inducción sobre la longitud de la secuencia, y también la normalización de un modelo de bigramas sin símbolo de fin. En el módulo~7, la expresión cerrada de las medias exponenciales de Adam se obtiene por inducción y se suma con la serie geométrica del \cref{ej:geometrica}. En el módulo~10, la corrección del modo directo de diferenciación automática se prueba por inducción sobre el orden topológico del grafo; en el módulo~12, la de BPE, por inducción sobre el número de fusiones; en el módulo~13, que la máscara causal impide que la salida en la posición $t$ dependa de posiciones posteriores; y en el módulo~17, la corrección del \emph{all-reduce} en anillo y de FlashAttention, por inducción sobre el número de pasos.
\end{observacion}

\begin{lecturas}
\item Hammack~\cite{hammack}, cap.~10 «Mathematical Induction» (§10.1–10.5).
\item Lehman, Leighton y Meyer~\cite{mcs}, cap.~2 «The Well Ordering Principle» (§2.1–2.3) y cap.~5 «Induction» (§5.1–5.3).
\item Velleman~\cite{velleman}, §6.1 «Proof by Mathematical Induction», §6.2 «More Examples» y §6.4 «Strong Induction».
\item MIT 6.042J~\cite{mit6042}, unidad~1, sesiones 3 (buen orden) y~8–9 (inducción).
\end{lecturas}

\section*{Ejercicios del tema 5}
\addcontentsline{toc}{section}{Ejercicios del tema 5}

\begin{ejercicio}{1}{eje:5-plantilla}
Para el enunciado «para todo $n\ge 1$, $1+3+5+\dots+(2n-1)=n^{2}$», escriba: (a) el predicado $P(n)$; (b) $P(1)$ y $P(n+1)$; (c) qué hay que demostrar exactamente en el paso de inducción. Después, complete la demostración.
\end{ejercicio}

\begin{ejercicio}{3}{eje:5-cuadrados}
Demuestre por inducción que, para todo $n\in\N$, $\displaystyle\sum_{k=1}^{n}k^{2}=\frac{n(n+1)(2n+1)}{6}$.
\fuente{adaptado de Hammack~\cite{hammack}, cap.~10 «Mathematical Induction».}
\end{ejercicio}

\begin{ejercicio}{3}{eje:5-cubos}
Demuestre por inducción que, para todo $n\in\N$, $\displaystyle\sum_{k=1}^{n}k^{3}=\Bigl(\frac{n(n+1)}{2}\Bigr)^{2}$.
\end{ejercicio}

\begin{ejercicio}{3}{eje:5-telescopica}
Demuestre por inducción que, para todo $n\in\N$, $\displaystyle\sum_{k=1}^{n}\frac{1}{k(k+1)}=\frac{n}{n+1}$. Dé después una segunda demostración sin inducción, usando $\frac{1}{k(k+1)}=\frac1k-\frac1{k+1}$.
\end{ejercicio}

\begin{ejercicio}{3}{eje:5-potencia}
Demuestre que $2^{n}>n^{2}$ para todo entero $n\ge 5$. ¿Para qué valores $n\in\{1,2,3,4\}$ es cierta la desigualdad?
\end{ejercicio}

\begin{ejercicio}{3}{eje:5-div9}
Demuestre que, para todo $n\in\N_0$, $9\divides 4^{n}+6n-1$.
\fuente{adaptado de Lehman, Leighton y Meyer~\cite{mcs}, §5.1 «Ordinary Induction» (problemas de la sección).}
\end{ejercicio}

\begin{ejercicio}{3}{eje:5-bernoulli}
Demuestre la desigualdad de Bernoulli: si $x\in\R$ y $x\ge -1$, entonces $(1+x)^{n}\ge 1+nx$ para todo $n\in\N_0$. Indique dónde se usa la hipótesis $x\ge -1$.
\end{ejercicio}

\begin{ejercicio}{3}{eje:5-sellos}
Demuestre por inducción fuerte que todo entero $n\ge 12$ se escribe como $n=4a+5b$ con $a,b\in\N_0$. ¿Qué enteros positivos menores que~12 no admiten esa forma?
\fuente{adaptado de Lehman, Leighton y Meyer~\cite{mcs}, §5.2 «Strong Induction».}
\end{ejercicio}

\begin{ejercicio}{3}{eje:5-fib-suma}
Sea $(F_n)$ la sucesión de Fibonacci ($F_1=F_2=1$, $F_{n+2}=F_{n+1}+F_n$). Demuestre que $\sum_{k=1}^{n}F_k=F_{n+2}-1$ para todo $n\in\N$.
\fuente{adaptado de Hammack~\cite{hammack}, §10.5 «Fibonacci Numbers».}
\end{ejercicio}

\begin{ejercicio}{3}{eje:5-fib-cota}
Demuestre por inducción fuerte que $F_n\le(5/3)^{n}$ para todo $n\in\N$. Explique por qué hacen falta dos casos base.
\fuente{adaptado de Hammack~\cite{hammack}, §10.5 «Fibonacci Numbers».}
\end{ejercicio}

\begin{ejercicio}{3}{eje:5-binario}
Demuestre por inducción fuerte que todo $n\in\N$ es suma de potencias de~2 distintas (es decir, tiene una representación binaria). (Indicación: distinga $n$ par y $n$ impar.)
\end{ejercicio}

\begin{ejercicio}{3}{eje:5-bo-geometrica}
Demuestre con el principio del buen orden (por el menor contraejemplo) que $\sum_{k=1}^{n}(2k-1)=n^{2}$ para todo $n\in\N$.
\fuente{adaptado de Lehman, Leighton y Meyer~\cite{mcs}, §2.2 «Template for Well Ordering Proofs».}
\end{ejercicio}

\begin{ejercicio}{2}{eje:5-caballos}
Se «demuestra» que, en todo conjunto de $n\ge 1$ caballos, todos tienen el mismo color: para $n=1$ es obvio; si vale para $n$, dado un conjunto de $n+1$ caballos se quita el primero y quedan $n$ del mismo color; se quita el último y quedan otros $n$ del mismo color; como ambos grupos comparten caballos, los $n+1$ tienen el mismo color. Señale con precisión el valor de $n$ para el que falla el paso y por qué.
\end{ejercicio}

\begin{ejercicio}{2}{eje:5-factorial}
(a) Calcule $S_n=\sum_{k=1}^{n}k\cdot k!$ para $n=1,\dots,5$ y conjeture una fórmula cerrada comparando con $(n+1)!$. (b) Demuestre la conjetura por inducción.
\end{ejercicio}

\begin{ejercicio}{2}{eje:5-sumas}
(a) Usando el \cref{ej:gauss}, obtenga una fórmula cerrada para $\sum_{k=1}^{n}(3k-1)$ y demuéstrela después por inducción. (b) Usando el \cref{ej:geometrica}, calcule $\sum_{k=0}^{7}2^{k}$ e interprete el resultado como el mayor número que cabe en un byte.
\end{ejercicio}

\begin{ejercicio}{3}{eje:5-log}
Sean $p_1,\dots,p_n$ números reales positivos. Demuestre por inducción sobre $n$, a partir de $\ln(ab)=\ln a+\ln b$, que $\ln\bigl(\prod_{t=1}^{n}p_t\bigr)=\sum_{t=1}^{n}\ln p_t$. Explique por qué esta identidad permite, en el módulo~1, convertir la probabilidad de una secuencia (un producto de muchos factores menores que~1) en una suma de log-probabilidades.
\end{ejercicio}
